\chapter{Analyticity, Crossing, and Landau Singularities}
\label{ch:volume-ii-analyticity-crossing-landau}
\ConventionsLorentzian


\section{Physical Region and First Sheet}
\label{sec:physical-region-and-first-sheet}

Consider \(2\to2\) scattering of identical stable scalar particles of mass
\(m\).  The metric is mostly plus, so a physical mass-\(m\) momentum \(k\)
obeys \(k^2=-m^2\).  Let the connected invariant amplitude be denoted by
\[
  \mathcal M(s,t),
\]
where
\[
  s=-(k_1+k_2)^2,\qquad
  t=-(k_1-k_3)^2,\qquad
  u=-(k_1-k_4)^2=4m^2-s-t .
\]
Here \(k_1,k_2\) are incoming and \(k_3,k_4\) are outgoing physical
momenta.  Equivalently, all four external momenta may be taken incoming after
replacing outgoing momenta by their negatives; the above definitions are the
physical scattering convention.  With \(k_i^2=-m^2\), the mostly-plus
Mandelstam identity is
\[
  s+t+u=4m^2 .
\]

In the center-of-mass frame,
\[
  s=E^2=4(k^2+m^2),
  \qquad
  t=-2k^2(1-\cos\theta),
  \qquad
  u=-2k^2(1+\cos\theta).
\]
Thus the physical \(s\)-channel region is
\[
  t\le0,\qquad
  u\le0,\qquad
  s\ge4m^2-t .
\]
The last inequality is the statement \(-t\le s-4m^2\).  For fixed
\(s\ge4m^2\), the allowed physical interval is
\[
  -(s-4m^2)\le t\le0 .
\]
For fixed real \(t<0\), the physical upper edge begins at
\[
  s=4m^2-t,
\]
because this is the value for which \(u=0\).  The analytic function whose
boundary value is the physical amplitude also has the lower \(s\)-channel
threshold at \(s=4m^2\); for \(4m^2<s<4m^2-t\) the point is on the same
right-hand cut but outside the physical \(s\)-channel scattering interval for
that fixed \(t\).

\begin{center}
\begin{tikzpicture}[>=Stealth,scale=0.92]
  \begin{scope}
    \node[draw,circle,fill=gray!15,minimum size=0.78cm,inner sep=0pt] (a) at (0,0) {};
    \draw[->,line width=0.9pt] (-1.75,0.75)--(a.150);
    \draw[->,line width=0.9pt] (-1.75,-0.75)--(a.210);
    \draw[->,line width=0.9pt] (a.30)--(1.75,0.75);
    \draw[->,line width=0.9pt] (a.330)--(1.75,-0.75);
    \node at (0,0) {\(\mathcal M\)};
    \draw[dashed] (0,0)--(1.45,0);
    \draw[blue!70!black,->,line width=0.9pt]
      (0.85,0) arc[start angle=0,end angle=24,radius=0.85];
    \node[blue!70!black] at (1.05,0.3) {\(\theta\)};
    \node[align=center] at (0,-1.55)
      {\(s=4(k^2+m^2)\)\\
       \(t=-2k^2(1-\cos\theta)\)};
  \end{scope}
  \begin{scope}[xshift=6.5cm]
    \draw[->,line width=0.75pt] (-0.3,0)--(5.2,0) node[right] {\(s\)};
    \draw[->,line width=0.75pt] (0,-1.45)--(0,1.85)
      node[above] {\(\operatorname{Im}s\)};
    \draw[red!75!black,decorate,
      decoration={snake,amplitude=0.55mm,segment length=2.7mm},
      line width=1.05pt] (2.05,0)--(4.95,0);
    \node[red!75!black,below] at (2.05,-0.15) {\(4m^2\)};
    \draw[orange!85!black,decorate,
      decoration={snake,amplitude=0.55mm,segment length=2.7mm},
      line width=1.05pt] (0.95,0)--(-0.1,0);
    \node[orange!85!black,below] at (0.95,-0.15) {\(-t\)};
    \fill[blue!75!black] (1.45,0.28) circle (2.1pt);
    \fill[blue!75!black] (1.68,0.28) circle (2.1pt);
    \node[blue!75!black,align=center,font=\small] at (1.55,0.85)
      {possible\\bound poles};
    \draw[purple!80!black,->,line width=1.0pt]
      (2.25,0.62) .. controls (3.05,1.18) and (3.82,0.92) .. (4.18,0.15);
    \node[purple!80!black,align=center,font=\small] at (3.52,1.18)
      {physical\\upper edge};
  \end{scope}
\end{tikzpicture}
\end{center}

For fixed real \(t<0\), the physical amplitude is the upper-edge boundary
value
\[
  \lim_{\epsilon\to0^+}\mathcal M(s+\ii\epsilon,t)
\]
on the first sheet.  In this monograph the terms \emph{first sheet} and
\emph{physical sheet} mean the same analytic branch: the germ reached from the
Euclidean or Feynman domain by the prescribed \(+\ii0\) continuation without
passing through any unitarity cut.  This is a statement about a branch of the
complexified invariant variables, not a synonym for the physical region of real
external momenta.  In the fixed-\(t\) \(s\)-plane this branch has the right-hand
\(s\)-channel cut
\[
  C_s=[4m^2,\infty),
\]
and, because \(u=4m^2-s-t\), the crossed \(u\)-channel cut
\[
  C_u=(-\infty,-t] .
\]
A stable state of mass \(M_B<2m\) coupled to the two external scalars gives a
first-sheet pole at \(s=M_B^2\) in the corresponding channel.  In the crossed
channel this same statement appears as a pole at
\[
  s=4m^2-t-M_B^2 .
\]
If the stable state is composite, this pole is the same spectral pole whose
four-point residue is factorized in
Theorem~\ref{thm:bethe-salpeter-pole-factorization}.  In a two-constituent
description its mass is equivalently determined by the homogeneous
Bethe--Salpeter condition
\eqref{eq:homogeneous-bethe-salpeter-equation}, or by
\((\mathbf 1-\widehat V_P)\Psi_B=0\) at \(P^2=-M_B^2\).  Crossing changes the
Mandelstam coordinate in which the pole is seen; it does not change the Hilbert
space state or the residue matrix element from which the pole term is built.
For the purposes used here, the first-sheet data consist of this cut plane
together with the possible pole data and with any additional first-sheet
singularities allowed by the full spectral problem.

A \emph{second sheet} is always defined relative to a specified cut.  It is the
adjacent analytic germ obtained by continuing the first-sheet germ through that
cut.  Near a single two-particle threshold, with
\[
  \rho(s)=\sqrt{1-\frac{s_{\mathrm{th}}}{s}}
\]
chosen so that \(\rho(s)>0\) on the physical upper edge \(s>s_{\mathrm{th}}\),
the continuation through the corresponding unitarity cut sends
\(\rho\mapsto-\rho\).  With several thresholds there is no cut-independent
meaning of ``the'' second sheet; the sheet is labelled by the set of cuts
through which the analytic germ has been continued, or equivalently by the
monodromy choices for the associated channel momenta.

\section{Real Analyticity and Crossing}

For the analytic discussion in this chapter assume a mass gap, Poincare
invariance, locality, unitarity, and the existence of stable one-particle
external states to which LSZ applies.  These assumptions specify the
structural input used here, while the maximal analytic domain remains an
additional problem.  Amplitudes first defined as physical boundary values are
continued as functions of complex kinematic invariants, and the same symbol
\(\mathcal M(s,t)\) denotes this analytic continuation on whichever domain is
under discussion.  Polynomial boundedness is a further high-energy hypothesis
and is introduced only in the dispersion-relation chapter.

Hermiticity of the theory and the Feynman boundary prescription imply real
analyticity on conjugate domains:
\[
  \mathcal M(s^*,t^*)=\mathcal M(s,t)^*
\]
where both sides are evaluated on corresponding branches.  On a real segment
not lying on a cut, this says that the first-sheet value is real up to pole
singularities.  On a cut, it relates the lower-edge and upper-edge boundary
values.  In perturbation theory the same property is visible graph by graph
when the couplings are real.

Crossing states that the \(s\)-, \(t\)-, and \(u\)-channel physical amplitudes
are boundary values of the same analytically continued four-point function in
different real regions, with the external representation conjugated when an
outgoing particle is moved to the incoming side or conversely.  On the real
\((s,t)\)-plane these regions are separated by thresholds and cuts.  In the
complexified kinematic domain they are connected by analytic continuation,
with the path and sheet specified by the singularities encountered.

\begin{center}
\begin{tikzpicture}[>=Stealth,scale=0.86]
  \draw[->,line width=0.75pt] (-3.7,0)--(3.85,0) node[right] {\(s\)};
  \draw[->,line width=0.75pt] (0,-3.0)--(0,3.15) node[above] {\(t\)};
  \coordinate (S) at (1.35,0);
  \coordinate (T) at (0,1.35);
  \node[below] at (S) {\(4m^2\)};
  \node[left] at (T) {\(4m^2\)};
  \draw[gray!55,line width=0.7pt] (-1.6,2.95)--(3.35,-2.0)
    node[below right] {\(u=0\)};
  \fill[blue!18] (S)--(3.35,0)--(3.35,-2.0)--cycle;
  \draw[blue!70!black,line width=1.35pt] (S)--(3.35,0);
  \draw[blue!70!black,line width=1.35pt] (S)--(3.35,-2.0);
  \node[blue!70!black,align=center,fill=white,inner sep=1.2pt]
    at (2.55,-0.75) {\(s\)-channel\\\(12\to34\)};
  \fill[green!18] (T)--(0,3.0)--(-1.65,3.0)--cycle;
  \draw[green!55!black,line width=1.35pt] (T)--(0,3.0);
  \draw[green!55!black,line width=1.35pt] (T)--(-1.65,3.0);
  \node[green!55!black,align=center,fill=white,inner sep=1.2pt]
    at (-0.72,2.25) {\(t\)-channel\\\(1\bar3\to\bar2 4\)};
  \fill[orange!18] (0,0)--(0,-3.0)--(-3.0,-3.0)--(-3.0,0)--cycle;
  \draw[orange!85!black,line width=1.35pt] (0,0)--(0,-3.0);
  \draw[orange!85!black,line width=1.35pt] (0,0)--(-3.0,0);
  \node[orange!85!black,align=center,fill=white,inner sep=1.2pt]
    at (-2.1,-1.15) {\(u\)-channel\\\(1\bar4\to3\bar2\)};
  \draw[purple!80!black,->,line width=1.05pt]
    (2.45,-0.18) .. controls (1.1,1.42) and (-0.55,1.76) .. (-1.05,2.02);
  \node[purple!80!black,align=center,fill=white,inner sep=1pt]
    at (0.45,1.18) {analytic\\continuation};
\end{tikzpicture}
\end{center}

\section{Unitarity on the Physical Cut}

Crossing and real analyticity constrain how different boundary values fit
together.  On the physical \(s\)-channel cut there is also the Hilbert-space
constraint from the unitary scattering operator.  Let \(S_\ell(s)\) denote the
angular-momentum-\(\ell\) partial-wave scattering eigenvalue in the
identical-scalar two-particle sector, with the normalization fixed earlier.
The elastic \(2\to2\) amplitude is
\[
  \mathcal M(s,t)
  =
  -16\pi\ii
  \sqrt{\frac{s}{s-4m^2}}
  \sum_{\ell=0}^{\infty}
  (2\ell+1)P_\ell(\cos\theta)\bigl(S_\ell(s)-1\bigr),
\]
where
\[
  \cos\theta=1+\frac{2t}{s-4m^2}.
\]
Unitarity of \(S\) implies, for \(s\ge4m^2\),
\[
  |S_\ell(s)|\le1 .
\]
If no inelastic channel is open, namely for
\[
  4m^2\le s\le M_{\mathrm{inel}}^2 ,
\]
the partial wave is a phase and
\[
  |S_\ell(s)|=1 .
\]
Here \(M_{\mathrm{inel}}\) is the invariant mass threshold of the lightest
state outside the two-particle elastic channel.
These inequalities are boundary conditions on the upper edge of the physical
cut; they are not by themselves an analytic continuation principle.  They will
combine with angular analyticity in the following chapter.

\section{Perturbative Discontinuities and Cut Rules}

The Hilbert-space unitarity statement also determines the discontinuity of
perturbative amplitudes across a physical cut.  Write
\[
  S=\mathbf 1+\ii T,
  \qquad
  \operatorname{Disc}_s\mathcal M
  =
  \mathcal M(s+\ii0,t)-\mathcal M(s-\ii0,t).
\]
From \(S^\dagger S=\mathbf 1\) one obtains
\[
  -\ii(T-T^\dagger)=T^\dagger T .
\]
Taking connected matrix elements and removing the common momentum-conserving
delta function gives, on the \(s\)-channel physical cut,
\begin{equation}
  \operatorname{Disc}_s\mathcal M_{\beta\alpha}
  =
  \ii
  \sum_X
  \int \dd\Phi_X\,
  \mathcal M_{X\alpha}\,
  \mathcal M_{\beta X}^{\,*}.
  \label{eq:cutkosky-unitarity-discontinuity}
\end{equation}
Here \(X\) runs over on-shell intermediate states with total momentum
\(P_\alpha\), and \(\dd\Phi_X\) is the Lorentz-invariant phase-space measure,
including the factor \(1/N!\) for \(N\) identical particles when the state sum is
over unordered particles.

\begin{theorem}[Cutkosky cut rule for a perturbative physical cut]
Consider a perturbative contribution to a connected amplitude whose
\(s\)-channel discontinuity is produced by putting a set \(C\) of internal
lines on shell so that the graph separates into a left subamplitude and a right
subamplitude.  With the amplitude normalization fixed by \(S=\mathbf 1+\ii T\),
the discontinuity contribution of this cut is obtained by replacing each cut
propagating line of mass \(m_i\) by the on-shell measure
\[
  \frac{\dd^D q_i}{(2\pi)^D}\,
  (2\pi)\theta(q_i^0)\delta(q_i^2+m_i^2),
\]
multiplying the left and right lower-order amplitudes, complex conjugating the
amplitude on the lower-edge side, and integrating over the common on-shell
phase space.  Equivalently, the graph contribution is the corresponding term in
\eqref{eq:cutkosky-unitarity-discontinuity}.
\end{theorem}

\begin{proof}
Expand \(T=\sum_r T^{(r)}\) in the chosen interaction order.  The order-\(N\)
part of \(-\ii(T-T^\dagger)=T^\dagger T\) is
\[
  -\ii\bigl(T^{(N)}-T^{(N)\dagger}\bigr)
  =
  \sum_{r=1}^{N-1}T^{(r)\dagger}T^{(N-r)} .
\]
Insert the spectral resolution of the identity on the physical Hilbert space
between the two factors on the right.  For a sector of \(n\) stable particles of
masses \(m_i\), this resolution contains
\[
  \frac1{n!}
  \prod_{i=1}^n
  \left[
    \int\frac{\dd^D q_i}{(2\pi)^{D-1}}\,
    \theta(q_i^0)\delta(q_i^2+m_i^2)
  \right]
  (2\pi)^D\delta^{(D)}\!\left(P_\alpha-\sum_i q_i\right).
\]
After the external LSZ residues and the overall momentum-conserving delta
function are removed, this is precisely the phase-space integral in
\eqref{eq:cutkosky-unitarity-discontinuity}.  In a Feynman-graph expansion the
same insertion cuts the internal lines that carry the intermediate on-shell
momenta and leaves the two lower-order amplitudes on the two sides of the cut.
\end{proof}

For the two-particle equal-mass cut, the phase-space factor appearing in
\eqref{eq:cutkosky-unitarity-discontinuity} is explicit.  Let \(P^2=-s\) and
\[
  \dd\Phi_2(P)
  =
  \int
  \frac{\dd^4 q}{(2\pi)^3}\theta(q^0)\delta(q^2+m^2)\,
  \frac{\dd^4 r}{(2\pi)^3}\theta(r^0)\delta(r^2+m^2)\,
  (2\pi)^4\delta^{(4)}(P-q-r).
\]
In the center-of-mass frame \(P=(\sqrt s,\vec0)\) this becomes
\[
  \dd\Phi_2(P)
  =
  \frac{\dd^3\vec q}{(2\pi)^2\,4\omega_{\vec q}^{\,2}}\,
  \delta(\sqrt s-2\omega_{\vec q}).
\]
Using \(\omega_{\vec q}=\sqrt{\vec q^{\,2}+m^2}\), the radial delta function
fixes
\[
  |\vec q|=\frac12\sqrt{s-4m^2},
\]
and hence
\begin{equation}
  \int\dd\Phi_2(P)
  =
  \frac1{8\pi}
  \sqrt{1-\frac{4m^2}{s}}\,
  \theta(s-4m^2).
  \label{eq:two-body-phase-space-four-d}
\end{equation}
For identical intermediate particles this is multiplied by \(1/2!\) if the
state sum is over unordered two-particle states.  Thus a one-loop bubble cut
has discontinuity equal to \(\ii\) times the product of the two tree
subamplitudes times \eqref{eq:two-body-phase-space-four-d}, with this symmetry
factor when appropriate.  The square-root branch point is the same threshold
that the Landau equations identify below.

\section{Crossing From Locality and Tube Analyticity}

The structural origin of crossing is the same analytic mechanism used for
Wightman functions.  For four scalar fields, the vacuum distribution
\[
  W_4(x_1,x_2,x_3,x_4)
  =
  \bra{\vac}
  \widehat\phi(x_1)\widehat\phi(x_2)
  \widehat\phi(x_3)\widehat\phi(x_4)
  \ket{\vac}
\]
has Fourier support fixed by the spectrum condition in relative variables.
It is therefore the boundary value of a holomorphic function in a tube
domain.  More explicitly, after removing the center-of-mass coordinate, the
Fourier transform has support in products of closed forward lightcones for
appropriate partial sums of momenta.  This support condition is the input
which makes the imaginary spacetime variables point into forward tubes.

Local commutativity says that, at real configurations where adjacent
insertions are spacelike separated,
\[
  W_4(\ldots,x_i,x_{i+1},\ldots)
  =
  W_4(\ldots,x_{i+1},x_i,\ldots)
\]
for identical scalar fields.  The edge-of-the-wedge theorem then identifies
the analytic continuations attached to the two orderings in the domain
generated from such spacelike configurations.

LSZ reduction is applied to the time-ordered or Wightman boundary values after
isolating the one-particle poles of the external legs.  Different assignments
of external legs as incoming or outgoing particles correspond to different
real boundary components of the same analytically continued four-point
function.  The \(s\)-channel process
\[
  1+2\to3+4
\]
and the crossed process
\[
  1+\bar{3}\to\bar{2}+4
\]
are therefore boundary values of one analytic object, with the Mandelstam
variables continued between their physical regions.  For a self-conjugate
scalar \(\bar{2}=2\) and \(\bar{3}=3\); for charged fields the bars indicate the
conjugate particle associated with the adjoint field.

The statement is local on the analytic domain.  It does not erase the
singularities separating physical regions.  Threshold cuts, bound-state poles,
and anomalous thresholds specify how the analytic continuation must be routed
and which boundary value is obtained at the end.  Crossing is the assertion of
common analytic origin; the singularity analysis below describes the
obstructions and boundary prescriptions encountered along particular
continuations.

\section{Edge-of-the-Wedge Form}

The analytic gluing used above has a precise local form.  Let \(F_+\) and
\(F_-\) be holomorphic functions in two domains
\[
  \mathcal D_+,\qquad \mathcal D_-
\]
of the complexified kinematic variables.  Suppose both domains have boundary
values on a common real open set \(U\), in the sense of distributional limits
against test functions supported in \(U\).  Suppose these boundary values
agree:
\[
  F_+|_U=F_-|_U .
\]
The edge-of-the-wedge theorem gives a holomorphic function \(F\) on a domain
containing \(\mathcal D_+\cup U\cup\mathcal D_-\) in its envelope of
holomorphy, with \(F=F_+\) on \(\mathcal D_+\) and \(F=F_-\) on
\(\mathcal D_-\).

In local QFT, the real set \(U\) is supplied by spacelike configurations where
local commutativity equates two Wightman orderings.  The domains
\(\mathcal D_\pm\) are the corresponding forward-tube domains selected by the
spectrum condition.  The theorem then produces a single analytic function
whose different real boundary components give differently ordered
correlators.  LSZ reduction transfers this analytic relation from correlators
to scattering amplitudes whenever the external one-particle pole hypotheses
hold.

\section{Pinches}

On the first sheet one finds stable one-particle poles, ordinary thresholds,
and anomalous thresholds.  An anomalous threshold is a contour-pinch
singularity on the physical sheet, caused by several internal lines going on
shell in a spacetime arrangement satisfying the positive Coleman--Norton
orientation condition defined below.  Its invariant location is determined by
the Landau configuration of those internal lines.

The singularities of perturbative amplitudes arise when an integration contour
is forced against singularities of the integrand.  Consider a scalar graph with
internal momenta
\[
  q_i=q_i(\ell_1,\ldots,\ell_L;p_1,\ldots,p_n),
\]
where the \(q_i\) are affine functions of the loop momenta \(\ell_a\) and
external momenta \(p_j\).  A typical denominator is
\[
  q_i^2+m_i^2-\ii0 .
\]
After Feynman parametrization, a product of \(N\) denominators has the form
\[
  \int_{\alpha_i\ge0}
  \prod_{i=1}^N\dd\alpha_i\,
  \delta\!\left(\sum_{i=1}^N\alpha_i-1\right)
  \frac{\text{numerator}}
       {\left[\sum_{i=1}^N\alpha_i(q_i^2+m_i^2)-\ii0\right]^N}.
\]
Numerator factors and ultraviolet subtractions can change the degree or even
remove a candidate singularity, but they do not change the necessary pinch
conditions for the scalar denominator.

As external momenta move in the complex domain, the loop contours can be
deformed as long as no singularity traps them.  A necessary condition for a
singularity is that the denominator vanish at a point where its first
variation in every loop-momentum direction also vanishes.  If the first
variation were nonzero, the contour could be displaced locally away from the
zero.  At a stationary zero, nearby singular hypersurfaces can pinch the
original real contour, and local deformations cannot remove the obstruction.

\begin{center}
\begin{tikzpicture}[>=Stealth,scale=0.95]
  \draw[blue!70!black,line width=1.25pt,->]
    (-3.3,-0.75) .. controls (-2.2,0.55) and (-0.6,0.95) .. (2.9,0.55);
  \fill[red!75!black] (-0.4,0.68) circle (2.5pt);
  \node[red!75!black,above] at (-0.4,0.78) {pole};
  \fill[orange!90!black] (-0.42,0.43) circle (2.5pt);
  \node[orange!90!black,left] at (-0.55,0.38) {contour};
  \draw[orange!90!black,->,line width=1.0pt] (-0.42,1.12)--(-0.42,0.55);
  \draw[orange!90!black,->,line width=1.0pt] (-0.42,-0.12)--(-0.42,0.30);
  \node[align=center] at (2.45,-0.9)
    {pinch: both local\\deformations are blocked};
\end{tikzpicture}
\end{center}

This stationary-pinch condition is the source of the Landau equations.  The
equations are necessary conditions for a singularity of a graph.  To decide
whether a singularity is actually present one must also check the sheet, the
orientation of the pinch, and whether the numerator cancels the local
singular term.

\section{Landau Equations}

Let
\[
  \mathcal D(\ell,\alpha;p)
  =
  \sum_{i=1}^N\alpha_i(q_i^2+m_i^2).
\]
A Landau singularity can occur only if there are Feynman parameters
\(\alpha_i\ge0\), not all zero, such that
\[
  \alpha_i(q_i^2+m_i^2)=0
  \qquad\text{for every internal line }i,
\]
and, for every loop momentum \(\ell_a^\mu\),
\[
  \frac{\partial\mathcal D}{\partial \ell_a^\mu}=0.
\]
Equivalently,
\[
  \sum_{i=1}^N
  \alpha_i q_{i\nu}
  \frac{\partial q_i^\nu}{\partial \ell_a^\mu}
  =
  0,
  \qquad a=1,\ldots,L .
\]
The omitted factor of \(2\) is irrelevant because the right hand side is
zero.

For a one-loop graph with all internal momenta oriented around the loop, this
reduces to
\[
  \sum_i\alpha_i q_i^\mu=0 .
\]

The first equation says that every line with \(\alpha_i>0\) is on shell:
\[
  q_i^2=-m_i^2 .
\]
The second equation says that the corresponding on-shell momenta form a
stationary configuration around each loop.  Solutions with some
\(\alpha_i=0\) describe singularities of reduced graphs in which the inactive
lines have been contracted.  Solutions with all relevant \(\alpha_i>0\) are
the leading Landau singularities of that graph.

With the terminology fixed in
Section~\ref{sec:physical-region-and-first-sheet}, the physical sheet is
selected by the sign of the Feynman prescription together with positive
Feynman parameters.  Positivity of the Feynman parameters distinguishes a
first-sheet threshold from the same algebraic locus reached only after
continuation through a cut.

\section{Spacetime Interpretation of Positive Landau Solutions}

The sign condition on the parameters has a spacetime meaning.  Consider a
Landau solution in a real physical region and keep only the active internal
lines, namely those with \(\alpha_i>0\).  Let the graph vertices be labelled
by \(v\), and orient each active edge \(i\) from a tail \(v_-(i)\) to a head
\(v_+(i)\).  The momentum \(q_i\) is the momentum assigned to this oriented
edge by the chosen loop-momentum routing.  The loop-stationarity equations say
that, for every oriented cycle \(C\) of the active graph,
\[
  \sum_{i\in C}\sigma_{Ci}\alpha_i q_i=0,\qquad
  \sigma_{Ci}=
  \begin{cases}
    +1,&\text{if the edge orientation agrees with the orientation of }C,\\
    -1,&\text{if it is opposite.}
  \end{cases}
\]
Equivalently, the vector assigned to an active edge by \(\alpha_i q_i\) has
zero sum around every closed path.  Hence, on each connected component, there
are vertex positions \(x_v\), unique up to a common translation, such that
\[
  x_{v_+(i)}-x_{v_-(i)}=\lambda_i q_i,
  \qquad
  \lambda_i\ge0 ,
\]
with \(\lambda_i\) proportional to \(\alpha_i\).  The proportionality constant
depends only on the convention used for the quadratic form in the parametrized
denominator.

The positive-energy condition is an additional orientation condition.  Choose a
time orientation in Minkowski space.  A positive Coleman--Norton realization of
the Landau solution consists of an orientation of every active edge in its
classical direction of propagation and momenta \(k_i\) on those oriented edges
such that
\[
  k_i^2=-m_i^2,\qquad k_i^0>0,\qquad
  x_{v_+(i)}-x_{v_-(i)}=\lambda_i k_i,\qquad
  \lambda_i\ge0 .
\]
If the loop-momentum routing originally assigned the opposite orientation to
that edge, then \(k_i=-q_i\) and the tail and head used in the spacetime
equation are interchanged.  Equivalently, one may reorient the active graph so
that \(q_i=k_i\) on every active edge.  Thus the compatibility condition is
precisely that each active internal momentum be a future-directed
positive-energy on-shell momentum in the direction in which the corresponding
edge is traversed in spacetime.

For \(m_i>0\), this is the worldline equation for a classical massive particle
propagating for nonnegative proper time, with \(\lambda_i m_i\) proportional to
that proper time.  For \(m_i=0\), it is the corresponding future-directed null
geodesic equation with nonnegative affine parameter.  Momentum conservation at
each vertex is the graph momentum conservation for the same oriented momenta,
with the external momenta assigned to that vertex.  Thus a real leading Landau
solution satisfying the positive Coleman--Norton orientation condition is a
classical spacetime process: all active internal particles go on shell, carry
positive energy in their direction of propagation, and meet at the vertices
with energy and momentum conserved.  This is the Coleman--Norton
interpretation.  It is a criterion for first-sheet physical singularities that
must be combined with the pinch analysis: a formal algebraic solution
contributes to the physical singularity only on the relevant sheet, with the
integration contour trapped, and with a numerator that leaves a nonzero local
singular term.

Subleading Landau solutions have some \(\alpha_i=0\).  Geometrically, the
corresponding internal lines are contracted, and the classical propagation
picture applies to the reduced graph.  Ordinary thresholds are reduced
classical processes in which a set of particles can be produced on shell with
the total external energy-momentum supplied to the subgraph.

\section{Two-Particle Threshold}

For the equal-mass two-propagator loop, take internal momenta
\[
  q_1=\ell,\qquad q_2=\ell+P .
\]
The Landau equations with \(0<\alpha<1\) are
\[
  \ell^2+m^2=0,\qquad
  (\ell+P)^2+m^2=0,
\]
and
\[
  (1-\alpha)\ell+\alpha(\ell+P)=0 .
\]
The stationary equation gives
\[
  \ell=-\alpha P .
\]
Substitution into the two on-shell equations gives
\[
  \alpha=\frac12,\qquad
  \ell=-\frac{P}{2},\qquad
  P^2=-4m^2 .
\]
Thus the ordinary two-particle threshold is a Landau singularity.
In terms of the invariant \(s=-P^2\), this is \(s=4m^2\).

\begin{center}
\begin{tikzpicture}[>=Stealth,scale=0.92]
  \begin{scope}
    \draw[line width=0.85pt] (-1.55,0.52)--(-0.48,0.52);
    \draw[line width=0.85pt] (-1.55,-0.52)--(-0.48,-0.52);
    \draw[line width=0.85pt] (0.48,0.52)--(1.55,0.52);
    \draw[line width=0.85pt] (0.48,-0.52)--(1.55,-0.52);
    \draw[line width=0.85pt] (0,0) ellipse (0.58 and 0.5);
    \node[blue!70!black] at (0,0.9) {\(q_2=\ell+P\)};
    \node[blue!70!black] at (0,-0.9) {\(q_1=\ell\)};
    \node at (0,-1.45) {bubble graph};
  \end{scope}
  \begin{scope}[xshift=5.3cm]
    \draw[->,line width=1.0pt] (0,0)--(-1.25,0) node[left] {\(q_1\)};
    \draw[->,line width=1.0pt] (0,0)--(1.25,0) node[right] {\(q_2\)};
    \draw[blue!70!black,->,line width=0.95pt] (1.25,-0.28)--(-1.25,-0.28)
      node[midway,below] {\(P\)};
    \node[align=center] at (0,0.85)
      {stationary point\\\(\alpha=\frac12\)};
    \node[blue!70!black] at (0,-1.05) {\(P^2=-4m^2\)};
  \end{scope}
\end{tikzpicture}
\end{center}

For unequal masses, the same equations give
\[
  \ell=-\alpha P,\qquad
  \alpha^2P^2=-m_1^2,\qquad
  (1-\alpha)^2P^2=-m_2^2 .
\]
The positive-parameter solution with \(0<\alpha<1\) is
\[
  \alpha=\frac{m_1}{m_1+m_2},
  \qquad
  P^2=-(m_1+m_2)^2 .
\]
This is the ordinary threshold: the two on-shell internal particles propagate
with zero relative spatial momentum in the center-of-mass frame.  The
algebraic equations also contain the pseudo-threshold
\[
  P^2=-(m_1-m_2)^2,
\]
but its parameter signs do not describe the same first-sheet physical
production process.

\section{Triangle Singularities}

The triangle graph is the first elementary example of a first-sheet singularity
produced by a genuine on-shell contour pinch rather than by a pole or endpoint
threshold. Let the three external momenta be
\[
  P_1^2=-M_1^2,\qquad
  P_2^2=-M_2^2,\qquad
  P_3^2=-M_3^2 ,
\]
with \(P_1+P_2+P_3=0\). The three internal lines are required to be on shell,
and the positive-parameter stationary condition requires their momentum
vectors to close consistently around the loop.
For one choice of loop orientation this means
\[
  q_i^2=-m_i^2,\qquad
  \alpha_1q_1+\alpha_2q_2+\alpha_3q_3=0,
  \qquad
  \alpha_i>0 .
\]
Together with momentum conservation at the vertices, these equations define
the candidate singular locus in the external invariants.  Equivalently, the
Gram matrix \(G_{ij}=q_i\cdot q_j\) has a null vector
\((\alpha_1,\alpha_2,\alpha_3)\) with strictly positive entries.  Thus the
condition is not only \(\det G=0\), but the existence of a positive null vector
for the oriented on-shell momenta.

\begin{center}
\begin{tikzpicture}[>=Stealth,scale=0.86]
  \begin{scope}
    \coordinate (a) at (-1.25,0.25);
    \coordinate (b) at (1.0,0.78);
    \coordinate (c) at (0.55,-1.05);
    \draw[line width=0.9pt] (a)--(b)--(c)--cycle;
    \foreach \p in {a,b,c}
      \fill[white,draw,line width=0.75pt] (\p) circle (0.15);
    \draw[blue!70!black,->,line width=1pt] (-2.35,0.25)--(a);
    \node[blue!70!black,above] at (-1.95,0.28) {\(P_1\)};
    \draw[blue!70!black,->,line width=1pt] (1.75,1.45)--(b);
    \node[blue!70!black,right] at (1.55,1.18) {\(P_3\)};
    \draw[blue!70!black,->,line width=1pt] (0.68,-2.15)--(c);
    \node[blue!70!black,right] at (0.62,-1.72) {\(P_2\)};
    \node at (0,-2.42) {triangle graph};
  \end{scope}
  \begin{scope}[xshift=5.35cm,yshift=-0.15cm]
    \draw[->,line width=1.0pt] (0,0)--(1.45,0);
    \node[right] at (1.53,0.08) {\(q_1\)};
    \draw[->,line width=1.0pt] (0,0)--(0.82,1.18);
    \node[above] at (0.78,1.35) {\(q_2\)};
    \draw[->,line width=1.0pt] (0,0)--(0.82,-1.18);
    \node[fill=white,inner sep=1pt] at (0.62,-0.74) {\(q_3\)};
    \draw[blue!70!black,->,line width=0.95pt] (0.82,1.18)--(1.82,0.38);
    \node[blue!70!black,right] at (1.9,0.55) {\(M_2\)};
    \draw[blue!70!black,->,line width=0.95pt] (0.82,-1.18)--(1.82,-0.38);
    \node[blue!70!black,right] at (1.9,-0.55) {\(M_1\)};
    \draw[blue!70!black,<->,line width=0.95pt] (0.22,-1.18)--(0.22,1.18);
    \node[blue!70!black,left] at (0.1,0) {\(M_3\)};
    \node[align=center] at (1.28,-2.12)
      {positive-parameter\\vector closure};
  \end{scope}
\end{tikzpicture}
\end{center}

In the equal-internal-mass planar example, the Euclideanized internal momentum
vectors in the sketch all have length \(m\).  The external invariant \(M_a\)
is the length of the chord between two such vectors, so
\[
  M_a^2=2m^2(1-\cos\varphi_a),
\]
where \(\varphi_a\) is the angle between the corresponding internal on-shell
vectors.  Thus \(M_a=\sqrt2\,m\) is the right-angle value.  If
\(M_1,M_2<\sqrt2\,m\), the relevant vectors lie in an open half-plane in the
configuration shown, and no positive relation
\(\alpha_1q_1+\alpha_2q_2+\alpha_3q_3=0\) exists.  If
\(M_1,M_2>\sqrt2\,m\), the vectors can be arranged with the origin in their
convex hull; the positive-parameter Landau equation can then be satisfied.
The remaining invariant is fixed by this geometry, giving a singularity at
the corresponding value \(P_3^2=-M_3^2\).

\begin{center}
\begin{tikzpicture}[>=Stealth,scale=0.84]
  \begin{scope}
    \node[align=center] at (0,1.95)
      {\(M_1,M_2<\sqrt2\,m\)};
    \draw[gray!55,dashed] (0,-1.45)--(0,1.45);
    \fill[gray!10] (0,-1.45)--(1.85,-1.45)--(1.85,1.45)--(0,1.45)--cycle;
    \coordinate (o) at (0,0);
    \coordinate (a) at (1.42,0.08);
    \coordinate (b) at (0.82,1.05);
    \coordinate (c) at (0.92,-0.72);
    \draw[green!50!black,dashed,line width=0.7pt] (a)--(b)--(c)--cycle;
    \draw[->,line width=1.0pt] (o)--(a) node[right] {\(q_1\)};
    \draw[->,line width=1.0pt] (o)--(b) node[above] {\(q_2\)};
    \draw[->,line width=1.0pt] (o)--(c) node[below] {\(q_3\)};
    \fill (o) circle (1.8pt);
    \node[align=center] at (0,-1.85)
      {origin outside convex hull\\no positive closure};
  \end{scope}
  \begin{scope}[xshift=6.4cm]
    \node[align=center] at (0,1.95)
      {\(M_1,M_2>\sqrt2\,m\)};
    \coordinate (o) at (0,0);
    \coordinate (a) at (1.28,0.0);
    \coordinate (b) at (-0.54,1.12);
    \coordinate (c) at (-0.66,-1.02);
    \fill[green!14] (a)--(b)--(c)--cycle;
    \draw[green!50!black,dashed,line width=0.75pt] (a)--(b)--(c)--cycle;
    \draw[->,line width=1.0pt] (o)--(a) node[right] {\(q_1\)};
    \draw[->,line width=1.0pt] (o)--(b) node[above left] {\(q_2\)};
    \draw[->,line width=1.0pt] (o)--(c) node[below left] {\(q_3\)};
    \fill (o) circle (1.8pt);
    \draw[blue!70!black,->,line width=0.9pt]
      (-0.1,0.25) .. controls (0.18,0.55) and (0.45,0.45) .. (0.62,0.2);
    \node[align=center] at (0,-1.85)
      {origin inside convex hull\\positive closure possible};
  \end{scope}
\end{tikzpicture}
\end{center}

The vector sketch represents the Coleman--Norton condition.  The actual
singular locus is obtained by solving the three on-shell equations together
with the positive-parameter stationary equation.  A solution gives a
first-sheet anomalous threshold at the determined value of \(P_3^2=-M_3^2\)
only when the active edges can be oriented with momenta \(k_i^0>0\) and
separations \(x_{v_+(i)}-x_{v_-(i)}=\lambda_i k_i\), \(\lambda_i\ge0\).  A
solution requiring reversed time ordering or a negative proper-time segment
lies on a continued sheet, even if it solves the same algebraic Landau
equations.

Such a singularity is an anomalous threshold: a first-sheet singularity of
the amplitude produced by on-shell internal propagation and contour pinching.
Its local data are the masses, momentum routing, and positive Feynman
parameters of the Landau solution.  Ordinary endpoint thresholds are the
reduced-graph case in which the on-shell process begins at the boundary of the
integration domain; the triangle is the first elementary interior-pinch case.
Coleman--Thun singularities in two-dimensional factorized scattering give a
nonperturbative setting in which the same on-shell propagation and contour
pinching mechanism appears.

\section{Use of Landau Analysis}

Landau equations are derived graph by graph. Their role here is diagnostic:
they identify the local analytic mechanisms by which thresholds and more
subtle first-sheet singularities arise in perturbation theory. In an exact
massive theory, the singularities of amplitudes are expected to be governed by
the same spectral and locality principles, with internal on-shell lines
replaced by actual stable particles or by spectral continua of the theory.

Within existing axiomatic frameworks, only parts of this analytic picture are
presently theorems. The classification used here is precise: stable particles
appear as first-sheet poles below thresholds, resonances as continued-sheet
poles, ordinary thresholds as multi-particle continua, and anomalous thresholds
as allowed on-shell propagation patterns encoded by the Landau equations.
This classification will be used in the following chapter only at the level
needed for angular analyticity: the nearest first-sheet crossed-channel
singularity fixes the distance from the physical angular interval to the edge
of the Lehmann ellipse.  Landau analysis supplies a concrete mechanism for
such singularities.  The nonperturbative construction of the \(S\)-operator
and the spectral assumptions stated earlier remain the structural input for
the exact scattering problem.
